\documentclass[10pt,letterpaper]{article}
\usepackage{cloudsuite}
\renewcommand{\DocShort}{CH. 14}
\begin{document}
\coverpage{Chapter 14: Advanced Cloud Architectures}{Cornell study notes and applied review | ARCHITECTURES}{Concepts \textbullet{} decisions \textbullet{} verification artifacts}
\sourceStatement
\tableofcontents
\newpage
\section{Learning contract}
\begin{takeaway}
Compare resilience approaches and enforce recovery/sovereignty constraints.
\end{takeaway}
\textbf{Chapter scope:} Clustering, continuity, recovery, sovereignty, provisioning and VPC.\\
\sourceRef{14}{371}
\subsection{Mapped source sections}
\begin{itemize}
\item 14.1 Hypervisor Clustering Architecture
\item 14.2 Virtual Server Clustering Architecture
\item 14.3 Load-Balanced Virtual Server Instances Architecture
\item 14.4 Nondisruptive Service Relocation Architecture
\item 14.5 Zero Downtime Architecture
\item 14.6 Cloud Balancing Architecture
\item 14.7 Resilient Disaster Recovery Architecture
\item 14.8 Distributed Data Sovereignty Architecture
\item 14.9 Resource Reservation Architecture
\item 14.10 Dynamic Failure Detection and Recovery Architecture
\item 14.11 Rapid Provisioning Architecture
\item 14.12 Storage Workload Management Architecture
\item 14.13 Virtual Private Cloud Architecture
\item 14.14 Case Study Example
\end{itemize}
\section{Cornell cue and notes matrix}
\small\setlength{\LTleft}{0pt}\setlength{\LTright}{0pt}
\begin{longtable}{>{\RaggedRight\arraybackslash}p{0.29\textwidth}>{\RaggedRight\arraybackslash}p{0.63\textwidth}}
\toprule\textbf{Cue / essential question}&\textbf{Explanatory notes and reasoning}\\\midrule\endfirsthead\toprule\textbf{Cue}&\textbf{Notes}\\\midrule\endhead
\textbf{How does hypervisor clustering differ from virtual server clustering?} & The former coordinates virtualization hosts; the latter coordinates workload guests. Their failure domains and maintenance mechanisms differ.\\[3pt]\textcolor{petrol}{\emph{Practice:} Test one host failure and one guest failure separately.}\\[8pt]\midrule
\textbf{What supports nondisruptive movement?} & Migration requires adequate destination capacity, network continuity and state transfer without violating service assumptions.\\[3pt]\textcolor{petrol}{\emph{Practice:} Measure errors and active sessions during relocation.}\\[8pt]\midrule
\textbf{How should zero-downtime claims be evaluated?} & Availability during upgrade or component loss depends on multiple mechanisms, including redundant instances, routing and compatibility.\\[3pt]\textcolor{petrol}{\emph{Practice:} Specify a rolling upgrade test with a loss budget.}\\[8pt]\midrule
\textbf{What distinguishes recovery from failover?} & Failover shifts active service after a fault; disaster recovery reconstructs service from broader disruptions and may accept explicit time or data loss.\\[3pt]\textcolor{petrol}{\emph{Practice:} Write RTO/RPO values with evidentiary tests.}\\[8pt]\midrule
\textbf{How do sovereignty and reservation constrain design?} & Jurisdiction and capacity guarantees can override the shortest or cheapest placement decision.\\[3pt]\textcolor{petrol}{\emph{Practice:} Compare two placements with compliance and capacity constraints.}\\[8pt]\midrule
\textbf{What can dynamic recovery and rapid provisioning accomplish?} & Automated detection/recovery and predescribed provisioning reduce manual delay, but require control over failure loops and configuration drift.\\[3pt]\textcolor{petrol}{\emph{Practice:} Test detection-to-repair time and rebuild reproducibility.}\\[8pt]\midrule
\textbf{Why isolate a virtual private cloud?} & A VPC architecture describes logical network segmentation inside a broader cloud environment; isolation is only as strong as its effective routing and policy.\\[3pt]\textcolor{petrol}{\emph{Practice:} Validate cross-segment connection denials.}\\[8pt]\midrule
\end{longtable}\normalsize
\section{Synthesis and architecture connections}
An architecture name is a repeatable composition of mechanisms with operating assumptions. Draw component relationships, failure domains, trust boundaries and a representative workload before committing to the pattern.
\begin{workbox}[title={Applied deliverable}]
Compare resilience approaches and enforce recovery/sovereignty constraints.\\[4pt]Create a one-page architecture or operations artifact with: (1) assumptions, (2) a decision and alternatives, (3) control ownership, (4) observable acceptance criteria, and (5) a failed-case test.
\end{workbox}
\section{Active recall}
\begin{enumerate}
\item How does hypervisor clustering differ from virtual server clustering?
\item What supports nondisruptive movement?
\item How should zero-downtime claims be evaluated?
\item What distinguishes recovery from failover?
\end{enumerate}
\subsection{Answer-check rubric}
\begin{itemize}
\item A correct response distinguishes the named concepts rather than relying on vendor names.
\item The explanation identifies an operational consequence, boundary or tradeoff.
\item At least one measurement, test or documentary artifact could demonstrate the claim.
\end{itemize}
\section{Blank Cornell capture sheet}
\textbf{Topic / source subsection:}\enspace\rule{0.65\linewidth}{0.4pt}\\[6pt]
\begin{longtable}{p{0.28\textwidth}p{0.64\textwidth}}
\toprule Cues / questions & Notes / evidence \\ \midrule
\rule{0pt}{42pt} & \\ \midrule
\rule{0pt}{42pt} & \\ \midrule
\rule{0pt}{42pt} & \\ \midrule
\rule{0pt}{42pt} & \\ \midrule
\rule{0pt}{42pt} & \\ \midrule
\end{longtable}
\textbf{Summary in three sentences:}\\[16pt]\hrulefill\\[16pt]\hrulefill\\[16pt]\hrulefill

\refs
\end{document}
